\section{Historian: where the ideas came from}
\label{sec:history}

The history of QC can be read as a change in what counts as a reference for
``normal'' data.

Early outlier work treated the problem as a statistical test. By the middle
of the twentieth century, Grubbs collected practical procedures for deciding
whether one or more observations were too extreme to belong to a normally
distributed sample \citep{grubbs1969}. The attraction was clear: the method
gave a named statistic, a significance level and a defensible rule.

The weakness was also clear. Classical tests use the sample mean and standard
deviation, both of which are affected by the very outliers they are meant to
detect. Pearson and Chandra Sekar described the masking problem, where two
similar outliers inflate the standard deviation enough to hide each other
\citep{pearson1936}. Rosner later developed the generalised ESD procedure to
handle several possible outliers without stopping at the first insignificant
step \citep{rosner1983}.

Robust statistics changed the reference. Instead of estimating the centre
and spread with quantities that a few bad values can move strongly, robust
methods use the median and the median absolute deviation (MAD). Iglewicz and
Hoaglin turned that idea into a simple engineering rule: compute a modified
z-score and flag large values, commonly using 3.5 as a screening threshold
\citep{iglewicz1993}. This was designed for batches of observations, not for
a short moving window over a correlated stream, which becomes important in
Section~\ref{sec:experiment}.

The next change was operational. The Oklahoma Mesonet made QC a system rather
than a single test. Its design combined calibration, maintenance, automatic
checks, graded flags and human review \citep{shafer2000}. The important idea
was not a new equation. It was that reliable data quality comes from several
independent kinds of evidence and a process for acting on the flags.

Campbell and colleagues generalised this approach for streaming
environmental sensors \citep{campbell2013}. They separate quality assurance
(what is done before data collection) from quality control (what is done to
the data afterwards), distinguish faults from real anomalies, and argue for
preserving raw values together with QC flags.

Leigh and colleagues then returned to explicit statistical comparison. On
high-frequency water-quality data they compared rule-based, regression-based
and feature-based methods and reported performance by fault type
\citep{leigh2019}. Their result is important for this report: no method wins
for every fault, so combinations are more useful than a search for one
universal detector.

\begin{figure}[h]
  \centering
  \includegraphics[width=0.93\textwidth]{fig_timeline.png}
  \caption{A compact timeline of the line followed in this report. The five
  core readings are marked in the labels.}
  \label{fig:timeline}
\end{figure}

The historical pattern is therefore simple. The score used to measure
surprise changed slowly. The reference for normality moved from theoretical
distributions, to robust summaries of the data, to engineering knowledge,
neighbouring stations and fitted models. The threshold remained the part
most likely to be inherited without enough justification.

% =====================================================================
% 3. Technician
% =====================================================================
